\documentclass[11pt]{article}

\usepackage[margin=0.82in]{geometry}
\usepackage{amsmath,amssymb}
\usepackage{array}
\usepackage{enumitem}
\usepackage{xcolor}
\usepackage{xurl}
\usepackage[hidelinks]{hyperref}

\setlength{\parindent}{0pt}
\setlength{\parskip}{0.45em}
\setlist[itemize]{topsep=0.2em,itemsep=0.15em,leftmargin=1.2em}
\setlist[enumerate]{topsep=0.2em,itemsep=0.35em,leftmargin=1.5em}

\newcommand{\blankline}{\vspace{0.9em}\hrule\vspace{0.9em}}
\newcommand{\shortblank}{\rule{0.35\linewidth}{0.4pt}}
\newcommand{\stepbox}[1]{%
  \vspace{0.25em}
  \fcolorbox{black!35}{black!3}{%
    \begin{minipage}{0.96\linewidth}
      #1
    \end{minipage}
  }
}

\begin{document}

\begin{center}
{\Large \textbf{Deriving the 1D Kalman Filter}}\\
\end{center}


\section*{1. Intuition: hidden state and noisy observations}

Suppose a sensor measures the position of an object on a line. The sensor is
noisy, so the raw readings $y_t$ jump around. The sensor noise variance is $R$.

\begin{enumerate}
  \item Suppose $R=4$. If you only think about the next raw observation $y_5$,
  what noise variance do you expect for this next sensor reading?

  \vspace{3em}

  \item Suppose we have strong reason to believe the object is stationary. Then
  previous readings can tell us where the fixed position is likely to be. Let
  $\mu_T$ be the average of the first $T$ readings:
  \[
  \mu_T = \frac{1}{T}\sum_{t=1}^T y_t.
  \]
  If the readings are independent and each has noise variance $R$, what is
  $P_T=\operatorname{Var}(\mu_T)$?

  (Note that this is different from the sample variance of the four readings)

  \vspace{3em}

  \item Suppose
  \[
  y_1=7,\qquad y_2=9,\qquad y_3=8,\qquad y_4=8.
  \]
  Compute the belief mean and variance after four readings, using $R=4$.
  \[
  \mu_4 =
  \hspace{5cm}
  P_4 =
  \hspace{5cm}
  \]

  \vspace{1.4em}

  \item A new reading arrives: $y_5=10$. Update the average without recomputing
  the whole sum:
  \[
  \mu_5
  =
  \mu_4+\alpha(y_5-\mu_4).
  \]
  What is $\alpha$? What is $\mu_5$?
  \[
  \alpha =
  \hspace{5cm}
  \mu_5 =
  \hspace{5cm}
  \]

  \vspace{1.4em}

  \item Suppose we have observed $N$ readings. We can compute the average
  $\mu_N$ and its variance. Now a new reading arrives: $y_{N+1}$. How should we
  update the average? How is the update weight $\alpha$ related to the variance
  of the average and the variance of the sensor noise?

  \vspace{3em}
\end{enumerate}

\newpage

\section*{2. Intuition: choosing the gain from uncertainty}
Now write the same idea as a belief update. Before seeing a new observation, suppose
\[
  x\sim \mathcal{N}(\mu_b,P_b),
  \qquad
  y\mid x\sim \mathcal{N}(x,R)
\]
A natural update is:
\[
  \mu_{\text{new}}=\mu_b+K(y-\mu_b) \qquad\qquad K=\frac{P_b}{P_b+R}
\]

\begin{enumerate}
  \item Check the limiting cases.
  \[
    R\to 0
    \quad\Rightarrow\quad
    K\to
    \hspace{3cm}
    R\to \infty
    \quad\Rightarrow\quad
    K\to
    \hspace{3cm}
  \]

  \vspace{1.4em}

  \[
    P_b\to 0
    \quad\Rightarrow\quad
    K\to
    \hspace{3cm}
    P_b\to \infty
    \quad\Rightarrow\quad
    K\to
    \hspace{3cm}
  \]

  \vspace{1.8em}

  \item Derive the variance update from the mean update.

  \vspace{5cm}

  \item Suppose:
  \[
  \mu_b = 8, \qquad P_b = 4, \qquad y = 10, \qquad R = 1.
  \]
  Is the new mean closer to $8$ or to $10$? Explain using uncertainty.

  \vspace{2.0em}

  Compute $K$, $\mu_{\text{new}}$, and $P_{\text{new}}$.
  \[
  K =
  \hspace{3cm}
  \mu_{\text{new}} =
  \hspace{3cm}
  P_{\text{new}} =
  \hspace{3cm}
  \]

  \vspace{2.0em}

\end{enumerate}

\newpage

\section*{3. Bayesian interpretation and the MAP solution}

Now we give the same update a Bayesian reading. The hidden position $x$ is not
observed directly. Before seeing the new measurement, our belief is the
\textbf{prior}:
\[
x \sim \mathcal{N}(\mu_b,P_b),
\qquad
p(x)\propto
\exp\left[-\frac{1}{2}\frac{(x-\mu_b)^2}{P_b}\right].
\]

The sensor model says that, if the true position were $x$, the measurement
would be centered at $x$ with noise variance $R$. This is the
\textbf{likelihood}:
\[
y\mid x \sim \mathcal{N}(x,R),
\qquad
p(y\mid x)\propto
\exp\left[-\frac{1}{2}\frac{(y-x)^2}{R}\right].
\]

Bayes' rule says to multiply the likelihood by the prior:
\[
p(x \mid y)
\propto
p(y\mid x)p(x)
\propto
\exp\left[
-\frac{1}{2}
\left(
\frac{(x-\mu_b)^2}{P_b}
+
\frac{(y-x)^2}{R}
\right)
\right].
\]

Define the quadratic cost:
\[
J(x)=
\frac{(x-\mu_b)^2}{P_b}
+
\frac{(y-x)^2}{R}.
\]
The \textbf{MAP(Maximum A Posteriori) estimate} is the value of $x$ that makes the posterior largest:
\[
x_{\mathrm{MAP}}=\arg\max_x p(x\mid y).
\]
Because $p(x\mid y)\propto \exp[-J(x)/2]$, maximizing the posterior is the same
as minimizing this quadratic cost:
\[
x_{\mathrm{MAP}}=\arg\min_x J(x).
\]

\begin{enumerate}
  \item Find the MAP estimate by differentiating $J(x)$ and setting the derivative to zero. With $K=P_b/(P_b+R)$, show that this is the same mean update as before:
  \[
  x_{\mathrm{MAP}}
  =
  \mu_b+K(y-\mu_b).
  \]

  \vspace{5.0cm}

  \item The posterior is Gaussian. For a Gaussian posterior, the MAP estimate
  and the posterior mean are the same, so $x_{\mathrm{MAP}}=\mu_{\text{new}}$.
  The variance comes from the curvature of $J(x)$:
  \[
  \frac{1}{P_{\text{new}}}
  =
  \hspace{5em}
  \qquad\Rightarrow\qquad
  P_{\text{new}}
  =
  \hspace{5em}
  \]

  Fill in the gain form:
  \[
  P_{\text{new}}
  =
  (1-K)P_b.
  \]

  \vspace{2.0em}
\end{enumerate}

\stepbox{
\textbf{Bayesian interpretation.}
For Gaussian belief and Gaussian sensor noise, the MAP solution gives exactly
the same gain update:
\[
K=\frac{P_b}{P_b+R},
\qquad
\mu_{\text{new}}=\mu_b+K(y-\mu_b),
\qquad
P_{\text{new}}=(1-K)P_b.
\]
The number $K$ tells how strongly the observation pulls the belief.
}

\newpage

\section*{4. Moving target: the 1D Kalman filter}

A Kalman filter repeats two steps: first move the belief forward with a model,
then update it using the new measurement. For a simple moving target:
\[
x_t = ax_{t-1}+e_t,
\qquad
e_t\sim\mathcal{N}(0,Q),
\qquad
y_t=x_t+v_t,
\qquad
v_t\sim\mathcal{N}(0,R).
\]

Before using $y_t$, the model gives:
\[
\mu_t^m=a\mu_{t-1},
\qquad
P_t^m=a^2P_{t-1}+Q.
\]

Then reuse the single-measurement update from Section 2:
\[
K_t=\frac{P_t^m}{P_t^m+R},
\qquad
\mu_t=\mu_t^m+K_t(y_t-\mu_t^m),
\qquad
P_t=(1-K_t)P_t^m.
\]

\textbf{Exercise.} Use
\[
\mu_0=0,\qquad P_0=2,\qquad a=1,\qquad Q=2,\qquad R=4,
\]
\[
y_1=2,\qquad y_2=2,\qquad y_3=5.
\]

For each row, first compute the model-based belief $(\mu_t^m,P_t^m)$, then
update it with the measurement $y_t$. Keep three decimal places if needed.

{\renewcommand{\arraystretch}{2.25}
\[
\begin{array}{c|c|c|c|c|c|c}
t & y_t & \mu_t^m & P_t^m & K_t & \mu_t & P_t \\
\hline
1 & 2 & \hspace{4.2em} & \hspace{4.2em} & \hspace{4.2em} &
\hspace{4.2em} & \hspace{4.2em} \\
2 & 2 & & & & & \\
3 & 5 & & & & &
\end{array}
\]
}

\vspace{2.0em}

Show the update for time step 1:
\[
\mu_1^m=
\qquad
P_1^m=
\qquad
K_1=
\]
\[
\mu_1=
\qquad
P_1=
\]

\vspace{2.5em}

\textbf{Interpret the result.}

\begin{enumerate}
  \item At time step 3, the sensor observation jumps to $5$. Does the belief
  jump all the way to $5$? Why or why not?

  \blankline
  \blankline

  \item Suppose the sensor becomes much noisier, so $R=16$ instead of $R=4$.
  Would the measurement weight become larger or smaller? What would that do to
  the new mean?

  \blankline
  \blankline

  \item In one sentence, answer the main filtering question: how does the filter
  reconcile disagreement between observation and belief?

  \blankline
  \blankline
\end{enumerate}

\newpage

\section*{5. Optional harder extension: two-dimensional position}

Now the hidden state is a vector:
\[
\mathbf{x}_t =
\begin{bmatrix}
x_t \\
y_t
\end{bmatrix}.
\]

Use the vector version of the moving-target model, where the sensor measures
both coordinates directly:
\[
\mathbf{x}_t=A\mathbf{x}_{t-1}+\mathbf{e}_t,
\qquad
\mathbf{z}_t=\mathbf{x}_t+\mathbf{v}_t.
\]

The two-dimensional equations are the vector version of the 1D belief update:
\[
\boldsymbol{\mu}_t^m = A\boldsymbol{\mu}_{t-1},
\qquad
P_t^m = AP_{t-1}A^\top+Q,
\]
\[
\mathbf{d}_t = \mathbf{z}_t-\boldsymbol{\mu}_t^m,
\qquad
S_t = P_t^m + R,
\qquad
K_t = P_t^m S_t^{-1},
\]
\[
\boldsymbol{\mu}_t =
\boldsymbol{\mu}_t^m + K_t\mathbf{d}_t,
\qquad
P_t = (I-K_t)P_t^m.
\]

Use:
\[
\boldsymbol{\mu}_0 =
\begin{bmatrix}
0 \\
0
\end{bmatrix},
\quad
P_0 =
\begin{bmatrix}
2 & 0 \\
0 & 2
\end{bmatrix},
\quad
A =
\begin{bmatrix}
1 & 0 \\
0 & 1
\end{bmatrix},
\]
\[
Q =
\begin{bmatrix}
1 & 0 \\
0 & 1
\end{bmatrix},
\quad
R =
\begin{bmatrix}
4 & 0 \\
0 & 1
\end{bmatrix},
\quad
\mathbf{z}_1 =
\begin{bmatrix}
2 \\
-1
\end{bmatrix}.
\]

\textbf{Important simplification:} all matrices here are diagonal. That means
the horizontal and vertical coordinates can be computed like two separate 1D
filters, but written in vector form.

\begin{enumerate}
  \item Carry the belief through the motion model.
  \[
  \boldsymbol{\mu}_1^m =
  \qquad\qquad
  P_1^m =
  \]

  \vspace{2.5em}

  \item Compute the discrepancy vector.
  \[
  \mathbf{d}_1 = \mathbf{z}_1-\boldsymbol{\mu}_1^m =
  \]

  \vspace{2.0em}

  \item Compute $S_1=P_1^m+R$ and then $S_1^{-1}$.
  \[
  S_1 =
  \qquad\qquad
  S_1^{-1} =
  \]

  \vspace{2.5em}

  \item Compute the measurement-weight matrix.
  \[
  K_1 = P_1^mS_1^{-1} =
  \]

  \vspace{2.5em}

  \item Compute the new mean and covariance.
  \[
  \boldsymbol{\mu}_1 =
  \boldsymbol{\mu}_1^m + K_1\mathbf{d}_1 =
  \]

  \vspace{2.0em}

  \[
  P_1 = (I-K_1)P_1^m =
  \]

  \vspace{2.5em}

  \item Which coordinate trusts the observation more: horizontal or vertical?
  How can you tell from $K_1$?

  \blankline
  \blankline
\end{enumerate}

\newpage

\section*{Reference notes used to design this exercise}

These references are useful for instructors who want to connect this scalar
exercise and the optional 2D extension to standard Kalman filter notation.

\begin{itemize}
  \item MIT, Tony Lacey, \emph{Tutorial: The Kalman Filter}.\\
  \url{https://web.mit.edu/kirtley/kirtley/binlustuff/literature/control/Kalman%20filter.pdf}

  \item Cornell University CS4758, \emph{Kalman Filter Applications: Tank Filling}.\\
  \url{https://www.cs.cornell.edu/courses/cs4758/2012sp/materials/MI63slides.pdf}

  \item Columbia University, Max Welling, \emph{The Kalman Filter}.\\
  \url{https://sites.stat.columbia.edu/liam/teaching/neurostat-spr12/papers/hmm/KF-welling-notes.pdf}

  \item Carnegie Mellon University 16-831, \emph{Kalman Filtering}.\\
  \url{https://www.cs.cmu.edu/~16831-f14/notes/F11/16831_lecture17_mklingen.pdf}

  \item University of Washington AMath 482/582, \emph{Lecture 26: Multivariate Kalman Filtering}.\\
  \url{https://atmos.washington.edu/~breth/classes/AS552/lect/lect26.pdf}
\end{itemize}

\end{document}
